\documentclass[a4paper,10pt]{article}
\usepackage[utf8]{inputenc}
\usepackage[T1]{fontenc}
\usepackage[french]{babel}
\usepackage{lmodern}
\usepackage[scaled=0.85]{beramono}
\usepackage[margin=2cm]{geometry}
\usepackage{xcolor}
\usepackage{listings}
\usepackage{textcomp}
\usepackage{fancyhdr}
\usepackage[hidelinks]{hyperref}

\definecolor{fond}{RGB}{248,248,248}
\definecolor{motcle}{RGB}{0,0,160}
\definecolor{commentaire}{RGB}{0,120,0}
\definecolor{chaine}{RGB}{163,21,21}
\definecolor{numero}{RGB}{140,140,140}

\lstdefinestyle{arduino}{
  language=C++,
  basicstyle=\ttfamily\footnotesize,
  keywordstyle=\color{motcle}\bfseries,
  commentstyle=\color{commentaire}\itshape,
  stringstyle=\color{chaine},
  numbers=left,
  numberstyle=\tiny\color{numero},
  numbersep=8pt,
  backgroundcolor=\color{fond},
  frame=single,
  rulecolor=\color{numero},
  breaklines=true,
  postbreak=\mbox{\textcolor{numero}{$\hookrightarrow$}\space},
  showstringspaces=false,
  tabsize=2,
  columns=fullflexible,
  keepspaces=true,
  upquote=true,
  extendedchars=true,
  morekeywords={uint8_t,uint16_t,uint32_t,size_t,byte,boolean,String,PROGMEM},
  literate={é}{{\'e}}1 {è}{{\`e}}1 {ê}{{\^e}}1 {à}{{\`a}}1 {ç}{{\c{c}}}1
           {ù}{{\`u}}1 {ô}{{\^o}}1 {î}{{\^i}}1 {É}{{\'E}}1 {·}{{\textperiodcentered}}1
}
\lstset{style=arduino}

\pagestyle{fancy}
\fancyhf{}
\fancyhead[L]{Snake ESP32 -- code source}
\fancyhead[R]{\nouppercase{\rightmark}}
\fancyfoot[C]{\thepage}

\title{Snake ESP32\\[0.3em]\large Code source du projet}
\author{}
\date{25/11/2026}

\begin{document}
\maketitle
\tableofcontents
\newpage
\section{Snake.ino}
Programme principal : setup() et boucle de jeu.

\begin{lstlisting}
#include <Arduino.h>
#include "config.h"
#include "etat.h"
#include "jeu.h"
#include "affichage.h"
#include "son.h"
#include "reseau.h"

void setup() {
  Serial.begin(115200);

  pinMode(SW_pin, INPUT_PULLUP);
  pinMode(BUZZER_pin, OUTPUT);

  
  initAffichage();
  initialiserJeu();
  initWifi();
}

void loop() {


  switch (etatJeux) {
    case MENU:
      mettreAJourMenu();
      break;

    case JEU:
      mettreAJourJeu();
      break;

    case GAME_OVER:
      mettreAJourGameOver();
      break;
  }

  switch (etatJeux) {
    case MENU:
      afficherMenu();
      break;

    case JEU:
      afficherJeu();
      break;

    case GAME_OVER:
      afficherGameOver();
      break;
  }

  mettreAJourWifi();
}
\end{lstlisting}

\section{config.h}
Constantes : broches, dimensions, vitesse, obstacles.

\begin{lstlisting}
#pragma once


const int SW_pin = 12;
const int X_pin = 34;
const int Y_pin = 35;
const int BUZZER_pin = 5;


const int LARGEUR_ECRAN = 128;
const int HAUTEUR_ECRAN = 64;
const int HAUT_ZONE_JEU = 16;  
const int grid = 8;


const int TAILLE_MAX_SERPENT = 20;   
const int LONGUEUR_MAX = 10;         


const int VITESSE_DEPART = 180;
const int VITESSE_MIN = 70;
const int PAS_VITESSE = 8;


const int MAX_OBS = 5;
const int FOOD_PAR_OBSTACLE = 3;


const int ZONE_MORTE = 500;
\end{lstlisting}

\section{etat.h}
Déclaration de l\textquotesingle{}état global du jeu.

\begin{lstlisting}
#pragma once
#include <Preferences.h>
#include "config.h"

enum Etat { MENU, JEU, GAME_OVER };


extern Preferences preferences;
extern Etat etatJeux;
extern int score;
extern int record;
extern int choixGameOver;


extern int snakeLength;
extern int snakeX[TAILLE_MAX_SERPENT];
extern int snakeY[TAILLE_MAX_SERPENT];
extern int dirX;
extern int dirY;
extern int nextDirX;
extern int nextDirY;


extern int foodX;
extern int foodY;
extern int obsX[MAX_OBS];
extern int obsY[MAX_OBS];
extern int nbObs;
extern int nbManges;


extern int gameSpeed;
extern unsigned long lastMove;
extern int centreX;
extern int centreY;
\end{lstlisting}

\section{etat.ino}
Définition de l\textquotesingle{}état global du jeu.

\begin{lstlisting}
#include "etat.h"


Preferences preferences;
Etat etatJeux = MENU;
int score = 0;
int record = 0;
int choixGameOver = 0;

int snakeLength = 3;
int snakeX[TAILLE_MAX_SERPENT];
int snakeY[TAILLE_MAX_SERPENT];
int dirX = 1;
int dirY = 0;
int nextDirX = 1;
int nextDirY = 0;

int foodX = 0;
int foodY = 0;
int obsX[MAX_OBS];
int obsY[MAX_OBS];
int nbObs = 0;
int nbManges = 0;

int gameSpeed = VITESSE_DEPART;
unsigned long lastMove = 0;
int centreX = 0;
int centreY = 0;
\end{lstlisting}

\section{jeu.h}
Interface de la logique de jeu.

\begin{lstlisting}
#pragma once

void initialiserJeu();       
void reset();               
void mettreAJourMenu();     
void mettreAJourJeu();     
void mettreAJourGameOver();
\end{lstlisting}

\section{jeu.ino}
Logique de jeu : serpent, nourriture, obstacles, collisions.

\begin{lstlisting}
#include <Arduino.h>
#include "jeu.h"
#include "config.h"
#include "etat.h"
#include "son.h"



static bool caseLibre(int x, int y) {
  for (int i = 0; i < snakeLength; i++) {
    if (snakeX[i] == x && snakeY[i] == y) return false;
  }
  for (int i = 0; i < nbObs; i++) {
    if (obsX[i] == x && obsY[i] == y) return false;
  }
  return true;
}

static void spawnFood() {
  do {
    foodX = random(0, LARGEUR_ECRAN / grid) * grid;
    foodY = random(2, HAUTEUR_ECRAN / grid) * grid;
  } while (!caseLibre(foodX, foodY));
}

static void spawnObstacle() {
  if (nbObs >= MAX_OBS) return;

  int x, y;
  do {
    x = random(0, LARGEUR_ECRAN / grid) * grid;
    y = random(2, HAUTEUR_ECRAN / grid) * grid;
  } while (!caseLibre(x, y)
           || (x == foodX && y == foodY)
           || (abs(x - snakeX[0]) <= 2 * grid && abs(y - snakeY[0]) <= 2 * grid));

  obsX[nbObs] = x;
  obsY[nbObs] = y;
  nbObs++;
}

void reset() {
  centreX = analogRead(X_pin);
  centreY = analogRead(Y_pin);

  snakeLength = 3;

  dirX = 1;
  dirY = 0;
  nextDirX = 1;
  nextDirY = 0;

  score = 0;
  gameSpeed = VITESSE_DEPART;
  nbObs = 0;
  nbManges = 0;

  snakeX[0] = 48;
  snakeY[0] = 32;

  snakeX[1] = 40;
  snakeY[1] = 32;

  snakeX[2] = 32;
  snakeY[2] = 32;

  spawnFood();
}

void initialiserJeu() {
  preferences.begin("snake", false);
  record = preferences.getInt("record", 0);
  reset();
}


static void lireJoystick() {


  int dx = analogRead(X_pin) - centreX;
  int dy = analogRead(Y_pin) - centreY;

  if (abs(dx) <= ZONE_MORTE && abs(dy) <= ZONE_MORTE) {
    return;
  }

  int nouvelleDirX = dirX;
  int nouvelleDirY = dirY;

  if (abs(dx) > abs(dy)) {
    nouvelleDirX = (dx > 0) ? 1 : -1;
    nouvelleDirY = 0;
  } else {
    nouvelleDirX = 0;
    nouvelleDirY = (dy > 0) ? 1 : -1;
  }

  if (nouvelleDirX != -dirX || nouvelleDirY != -dirY) {
    dirX = nouvelleDirX;
    dirY = nouvelleDirY;
  }
}



static void deplacerSerpent() {

  for (int i = snakeLength; i > 0; i--) {

    snakeX[i] = snakeX[i - 1];
    snakeY[i] = snakeY[i - 1];
  }

  snakeX[0] += dirX * grid;
  snakeY[0] += dirY * grid;
}

static void declencherGameOver() {

  if (score > record) {

    record = score;
    preferences.putInt("record", record);
  }

  choixGameOver = 0;
  etatJeux = GAME_OVER;

  bip(400);
}

void mettreAJourMenu() {

  if (digitalRead(SW_pin) == 0) {

    etatJeux = JEU;
    reset();
    bip(40);
    delay(200);
  }
}

void mettreAJourJeu() {

  if (millis() - lastMove <= gameSpeed) return;

  lastMove = millis();

  lireJoystick();
  deplacerSerpent();

  if (snakeX[0] < 0) {

    snakeX[0] = LARGEUR_ECRAN - grid;

  } 
  else if (snakeX[0] >= LARGEUR_ECRAN) 
  {
    snakeX[0] = 0;

  }

  if (snakeY[0] < HAUT_ZONE_JEU) {

    snakeY[0] = HAUTEUR_ECRAN - grid;

  } 
  else if (snakeY[0] >= HAUTEUR_ECRAN) {

    snakeY[0] = HAUT_ZONE_JEU;

  }


  for (int i = 0; i < nbObs; i++) {
    if (snakeX[0] == obsX[i] && snakeY[0] == obsY[i]) {
      declencherGameOver();
      return;
    }
  }


  for (int i = 1; i < snakeLength; i++) {
    if (snakeX[0] == snakeX[i] && snakeY[0] == snakeY[i]) {
      declencherGameOver();
      return;
    }
  }


  if (snakeX[0] == foodX && snakeY[0] == foodY) {

    if (snakeLength < LONGUEUR_MAX) snakeLength++;

    score += 2;
    nbManges++;

    gameSpeed -= PAS_VITESSE;
    if (gameSpeed < VITESSE_MIN) gameSpeed = VITESSE_MIN;

    bip(60);
    spawnFood();

    if (nbManges % FOOD_PAR_OBSTACLE == 0) {
      spawnObstacle();
    }
  }
}

void mettreAJourGameOver() {

  int x = analogRead(X_pin);
  if (x < 1000) {
    choixGameOver = 0;
  } else if (x > 3000) {
    choixGameOver = 1;
  }

  if (digitalRead(SW_pin) == 0) {
    bip(40);

    while (digitalRead(SW_pin) == 0) {
      delay(10);
    }

    etatJeux = (choixGameOver == 0) ? JEU : MENU;
    reset();
  }
}
\end{lstlisting}

\section{affichage.h}
Interface de l\textquotesingle{}affichage OLED.

\begin{lstlisting}
#pragma once

void initAffichage();
void afficherMenu();
void afficherJeu();
void afficherGameOver();
\end{lstlisting}

\section{affichage.ino}
Affichage OLED : menu, jeu, game over.

\begin{lstlisting}
#include <Arduino.h>
#include <U8g2lib.h>
#include "affichage.h"
#include "config.h"
#include "etat.h"


static U8G2_SSD1306_128X64_NONAME_F_HW_I2C u8g2(U8G2_R0, U8X8_PIN_NONE);

void initAffichage() {
  u8g2.setI2CAddress(0x3C << 1);
  u8g2.begin();
}
static void drawCentre(const char* txt, int y) {
  int x = (LARGEUR_ECRAN - u8g2.getStrWidth(txt)) / 2;
  u8g2.drawStr(x, y, txt);
}

void afficherMenu() {
  u8g2.clearBuffer();

  u8g2.setFont(u8g2_font_courB18_tf);
  drawCentre("SNAKE", 20);

  u8g2.setFont(u8g2_font_7x13_tf);
  String txt = "Record : " + String(record);
  drawCentre(txt.c_str(), 42);
  drawCentre("Clic pour jouer", 60);

  u8g2.sendBuffer();
}

void afficherJeu() {
  u8g2.clearBuffer();


  for (int i = 0; i < snakeLength; i++) {
    u8g2.drawBox(snakeX[i], snakeY[i], grid, grid);
  }


  u8g2.setFont(u8g2_font_6x10_tf);
  String txtScore = "Score: " + String(score);
  u8g2.drawStr(4, 10, txtScore.c_str());
  u8g2.drawHLine(0, 15, LARGEUR_ECRAN);


  for (int y = HAUT_ZONE_JEU; y < HAUTEUR_ECRAN; y += grid) {
    for (int x = 0; x < LARGEUR_ECRAN; x += grid) {
      u8g2.drawPixel(x, y);
    }
  }


  for (int i = 0; i < nbObs; i++) {
    u8g2.drawFrame(obsX[i], obsY[i], grid, grid);
    u8g2.drawLine(obsX[i], obsY[i], obsX[i] + grid - 1, obsY[i] + grid - 1);
    u8g2.drawLine(obsX[i] + grid - 1, obsY[i], obsX[i], obsY[i] + grid - 1);
  }


  u8g2.drawFrame(foodX, foodY, grid, grid);

  u8g2.sendBuffer();
}

void afficherGameOver() {
  u8g2.clearBuffer();

  u8g2.setFont(u8g2_font_courB12_tf);
  drawCentre("GAME OVER", 14);

  u8g2.setFont(u8g2_font_6x10_tf);
  String txtScore = "Score : " + String(score);
  drawCentre(txtScore.c_str(), 30);
  String txtRecord = "Record : " + String(record);
  drawCentre(txtRecord.c_str(), 42);

  u8g2.drawStr(14, 59, "Rejouer");
  u8g2.drawStr(84, 59, "Menu");

  if (choixGameOver == 0) {
    u8g2.drawFrame(10, 49, u8g2.getStrWidth("Rejouer") + 8, 13);
  } else {
    u8g2.drawFrame(80, 49, u8g2.getStrWidth("Menu") + 8, 13);
  }

  u8g2.sendBuffer();
}
\end{lstlisting}

\section{son.h}
Interface du buzzer.

\begin{lstlisting}
#pragma once

void bip(int ms);
\end{lstlisting}

\section{son.ino}
Gestion du buzzer.

\begin{lstlisting}
#include <Arduino.h>
#include "son.h"
#include "config.h"

void bip(int ms) {
  digitalWrite(BUZZER_pin, HIGH);
  delay(ms);
  digitalWrite(BUZZER_pin, LOW);
}
\end{lstlisting}

\section{reseau.h}
Interface réseau.

\textit{Fichier vide.}

\section{wifi.ino}
Wi-Fi, serveur web et WebSocket.

\begin{lstlisting}
#include <Arduino.h>
#include <WiFi.h>
#include <WebServer.h>
#include <WebSocketsServer.h>

#include "reseau.h"
#include "etat.h"
#include "page.h"

const char* ssid = "ESP32-MonReseau";
const char* password = "12345678";

WebServer server(80);
WebSocketsServer webSocket(81);

String dernierEtatEnvoye;

String construireEtatJeu() {
  String json;
  json.reserve(1500);

  json = "{\"score\":" + String(score);
  json += ",\"record\":" + String(record);
  json += ",\"etat\":" + String((int)etatJeux);
  json += ",\"choixGameOver\":" + String(choixGameOver);
  json += ",\"foodX\":" + String(foodX);
  json += ",\"foodY\":" + String(foodY);
  json += ",\"serpent\":[";

  for (int i = 0; i < snakeLength; i++) {
    if (i > 0) {
      json += ",";
    }

    json += "{\"x\":" + String(snakeX[i]);
    json += ",\"y\":" + String(snakeY[i]) + "}";
  }

  json += "],\"obstacles\":[";

  for (int i = 0; i < nbObs; i++) {
    if (i > 0) {
      json += ",";
    }

    json += "{\"x\":" + String(obsX[i]);
    json += ",\"y\":" + String(obsY[i]) + "}";
  }

  json += "]}";

  return json;
}

void handleRoot() {
  server.sendHeader("Cache-Control", "no-store");
  server.send_P(200, "text/html; charset=utf-8", PAGE_HTML);
}

void initWifi() {
  WiFi.mode(WIFI_AP);

  if (!WiFi.softAP(ssid, password)) {
    Serial.println("Erreur : creation du reseau Wi-Fi impossible");
    return;
  }

  Serial.print("Reseau Wi-Fi : ");
  Serial.println(ssid);

  Serial.print("Adresse IP : ");
  Serial.println(WiFi.softAPIP());

  server.on("/", handleRoot);
  server.begin();

  webSocket.begin();

  webSocket.onEvent(
    [](uint8_t client, WStype_t type, uint8_t* payload, size_t length) {
      if (type == WStype_CONNECTED) {
        String json = construireEtatJeu();
        webSocket.sendTXT(client, json);
      }
    }
  );

  Serial.println("Serveur web et WebSocket demarres");
}

void mettreAJourWifi() {
  server.handleClient();
  webSocket.loop();

  String json = construireEtatJeu();

  if (json != dernierEtatEnvoye) {
    webSocket.broadcastTXT(json);
    dernierEtatEnvoye = json;
  }
}
\end{lstlisting}

\section{page.h}
Page web embarquée (HTML / CSS / JavaScript).

\begin{lstlisting}
#pragma once
#include <Arduino.h>

const char PAGE_HTML[] PROGMEM = R"rawliteral(
<!DOCTYPE html>
<html lang="fr">
<head>
  <meta charset="UTF-8">
  <meta name="viewport" content="width=device-width, initial-scale=1">
  <title>Snake ESP32</title>

  <style>
    * {
      box-sizing: border-box;
    }

    body {
      margin: 0;
      min-height: 100vh;
      padding: 24px;
      display: flex;
      align-items: center;
      justify-content: center;
      background: #0b121a;
      color: #eef6ff;
      font-family: Arial, sans-serif;
    }

    main {
      width: min(100%, 960px);
    }

    header, .scores, footer {
      display: flex;
      align-items: center;
      justify-content: space-between;
      gap: 16px;
      flex-wrap: wrap;
    }

    header {
      margin-bottom: 24px;
    }

    h1 {
      margin: 0;
      font-size: clamp(24px, 5vw, 40px);
    }

    #connexion {
      color: #ffca75;
      font-size: 14px;
    }

    .scores {
      margin-bottom: 16px;
      padding: 16px 20px;
      background: #17232f;
      border-radius: 12px;
      font-size: clamp(16px, 3vw, 24px);
    }

    strong {
      color: #a9ff83;
    }

    .terrain {
      position: relative;
      overflow: hidden;
      border: 2px solid #304555;
      border-radius: 14px;
      background: #111e29;
    }

    canvas {
      display: block;
      width: 100%;
      height: auto;
      image-rendering: pixelated;
    }

    .ecran {
      position: absolute;
      inset: 0;
      display: flex;
      flex-direction: column;
      align-items: center;
      justify-content: center;
      gap: 14px;
      padding: 12px;
      background: rgba(11, 18, 26, 0.95);
      text-align: center;
    }

    [hidden] {
      display: none !important;
    }

    h2 {
      margin: 0;
      color: #a9ff83;
      font-size: clamp(26px, 6vw, 64px);
    }

    p {
      margin: 0;
      color: #b2c4d2;
      font-size: clamp(12px, 2.2vw, 20px);
    }

    #gameOver h2 {
      color: #ff7589;
    }

    .choix {
      display: flex;
      gap: 16px;
    }

    .option {
      padding: 10px 24px;
      border: 2px solid #425768;
      border-radius: 10px;
      color: #b2c4d2;
      font-size: clamp(14px, 2.5vw, 22px);
    }

    .option.selection {
      color: #0b121a;
      background: #a9ff83;
      border-color: #a9ff83;
      font-weight: bold;
    }

    footer {
      margin-top: 18px;
      color: #91a8ba;
      font-size: 14px;
    }

    button {
      padding: 10px 16px;
      border: 1px solid #425768;
      border-radius: 8px;
      background: #17232f;
      color: #eef6ff;
      cursor: pointer;
    }

    @media (max-width: 550px) {
      body {
        padding: 12px;
      }

      .ecran {
        gap: 6px;
      }

      .option {
        padding: 6px 16px;
      }
    }
  </style>
</head>

<body>
  <main>
    <header>
      <h1>Snake ESP32</h1>
      <span id="connexion" role="status">Connexion...</span>
    </header>

    <div class="scores">
      <span>Score : <strong id="score">0</strong></span>
      <span>Record : <strong id="record">0</strong></span>
    </div>

    <div class="terrain">
      <canvas id="terrain" width="128" height="48"></canvas>

      <section id="attente" class="ecran">
        <h2>Connexion...</h2>
        <p>En attente du jeu sur l'ESP32</p>
      </section>

      <section id="menu" class="ecran" hidden>
        <h2>SNAKE</h2>
        <p>Record : <strong id="recordMenu">0</strong></p>
        <p>Appuie sur le joystick pour jouer</p>
      </section>

      <section id="gameOver" class="ecran" hidden>
        <h2>GAME OVER</h2>

        <p>
          Score : <strong id="scoreFinal">0</strong>
          | Record : <strong id="recordFinal">0</strong>
        </p>

        <div class="choix">
          <span id="rejouer" class="option">Rejouer</span>
          <span id="retourMenu" class="option">Menu</span>
        </div>

        <p>Gauche / droite pour choisir · Clic pour valider</p>
      </section>
    </div>

    <footer>
      <span>Commandes sur le joystick de l'ESP32</span>
      <button id="pleinEcran" type="button">Plein écran</button>
    </footer>
  </main>

  <script>
    const MENU = 0;
    const JEU = 1;
    const GAME_OVER = 2;

    const GRILLE = 8;
    const HAUT_ZONE_JEU = 16;

    const elements = {};

    [
      "terrain", "connexion", "attente", "menu", "gameOver",
      "score", "record", "recordMenu", "scoreFinal", "recordFinal",
      "rejouer", "retourMenu", "pleinEcran"
    ].forEach(id => {
      elements[id] = document.getElementById(id);
    });

    const canvas = elements.terrain;
    const ctx = canvas.getContext("2d");

    let socket = null;
    let reconnexion = null;

    function dessinerCase(x, y, couleur) {
      ctx.fillStyle = couleur;
      ctx.fillRect(
        x,
        y - HAUT_ZONE_JEU,
        GRILLE - 1,
        GRILLE - 1
      );
    }

    function dessinerTerrain(jeu) {
      ctx.fillStyle = "#111e29";
      ctx.fillRect(0, 0, canvas.width, canvas.height);

      ctx.fillStyle = "#2b3d4c";

      for (let y = 0; y < canvas.height; y += GRILLE) {
        for (let x = 0; x < canvas.width; x += GRILLE) {
          ctx.fillRect(x, y, 1, 1);
        }
      }

      for (const obstacle of jeu.obstacles) {
        dessinerCase(obstacle.x, obstacle.y, "#8295a3");

        ctx.fillStyle = "#354958";
        ctx.fillRect(
          obstacle.x + 2,
          obstacle.y - HAUT_ZONE_JEU + 2,
          3,
          3
        );
      }

      dessinerCase(jeu.foodX, jeu.foodY, "#ff6577");

      jeu.serpent.forEach((segment, index) => {
        dessinerCase(
          segment.x,
          segment.y,
          index === 0 ? "#b4ff75" : "#54ce87"
        );
      });
    }

    function afficherEtat(jeu) {
      if (
        ![MENU, JEU, GAME_OVER].includes(jeu.etat) ||
        !Array.isArray(jeu.serpent) ||
        !Array.isArray(jeu.obstacles)
      ) {
        throw new Error("Etat du jeu incomplet");
      }

      elements.score.textContent = jeu.score;
      elements.record.textContent = jeu.record;
      elements.recordMenu.textContent = jeu.record;
      elements.scoreFinal.textContent = jeu.score;
      elements.recordFinal.textContent = jeu.record;

      dessinerTerrain(jeu);

      elements.attente.hidden = true;
      elements.menu.hidden = jeu.etat !== MENU;
      elements.gameOver.hidden = jeu.etat !== GAME_OVER;

      elements.rejouer.classList.toggle(
        "selection",
        jeu.choixGameOver === 0
      );

      elements.retourMenu.classList.toggle(
        "selection",
        jeu.choixGameOver === 1
      );

      const noms = ["Menu", "Partie en cours", "Game Over"];

      elements.connexion.textContent = "Connecté · " + noms[jeu.etat];
      elements.connexion.style.color = "#a9ff83";
    }

    function connecterWebSocket() {
      if (
        socket &&
        (socket.readyState === WebSocket.OPEN ||
         socket.readyState === WebSocket.CONNECTING)
      ) {
        return;
      }

      clearTimeout(reconnexion);

      elements.connexion.textContent = "Connexion...";
      elements.connexion.style.color = "#ffca75";

      socket = new WebSocket(
        "ws://" + window.location.hostname + ":81/"
      );

      const connexionActuelle = socket;

      connexionActuelle.onopen = () => {
        elements.connexion.textContent = "Connecté · En attente du jeu";
        elements.connexion.style.color = "#a9ff83";
      };

      connexionActuelle.onmessage = (message) => {
        try {
          afficherEtat(JSON.parse(message.data));
        } catch (erreur) {
          elements.connexion.textContent = "Données du jeu invalides";
          elements.connexion.style.color = "#ffca75";
          console.error(erreur);
        }
      };

      connexionActuelle.onclose = () => {
        if (socket !== connexionActuelle) {
          return;
        }

        socket = null;
        elements.connexion.textContent =
          "Déconnecté · Nouvelle tentative...";
        elements.connexion.style.color = "#ffca75";

        reconnexion = setTimeout(connecterWebSocket, 1000);
      };

      connexionActuelle.onerror = () => {
        connexionActuelle.close();
      };
    }

    if (!document.fullscreenEnabled) {
      elements.pleinEcran.hidden = true;
    }

    elements.pleinEcran.addEventListener("click", async () => {
      try {
        if (document.fullscreenElement) {
          await document.exitFullscreen();
        } else {
          await document.documentElement.requestFullscreen();
        }
      } catch (erreur) {
        console.error(erreur);
      }
    });

    document.addEventListener("fullscreenchange", () => {
      elements.pleinEcran.textContent = document.fullscreenElement
        ? "Quitter le plein écran"
        : "Plein écran";
    });

    connecterWebSocket();
  </script>
</body>
</html>
)rawliteral";
\end{lstlisting}

\end{document}
